\section{Results}
\label{sec:results}

\subsection{Model Comparison}

We evaluate all four of our models on the same complete-case test set, which runs from January~2023 to March~2026. After cleaning, we are left with 36 test observations in San~Francisco and Austin and 25 in Cleveland, applied uniformly to all four models. The smaller Cleveland count reflects the fact that the extracted Cleveland LAUS metro unemployment series ends in December~2024, while the Zillow and national FRED series extend through March~2026; once the one-month lag is applied to \texttt{d\_unrate\_metro}, Cleveland's last complete-case test month is January~2025. We report the two metrics defined in Section~\ref{sec:method}: the mean absolute error in log returns and the dollar-denominated MAE.

\begin{table}[!htb]
\centering
\caption{Dollar MAE by model and metro (USD per month, test period: 2023-01 to 2026-03). Bold indicates the best performance for each metro in the row.}
\label{tab:dollar_mae}
\begin{tabular}{lrrrr}
\hline
Metro & OLS & ARIMA(AIC) & XGBoost & MLP(16,8) \\
\hline
San Francisco, CA & \textbf{\$3{,}898} & \$5{,}660 & \$4{,}210 & \$5{,}219 \\
Austin, TX        & \$2{,}574 & \$2{,}371 & \$3{,}494 & \textbf{\$2{,}286} \\
Cleveland, OH     & \$748     & \textbf{\$491} & \$549 & \$1{,}064 \\
\hline
\end{tabular}
\end{table}

Figure~\ref{fig:maebar} provides the dollar MAE results for
all model-metro pairings.

\begin{figure}[!htb]
\centering
\includegraphics[width=0.68\textwidth]{mae_dollars_bar.png}
\caption{Dollar MAE for metro areas and models (test period: 
2023-01 to 2026-03). The lower bars represent better forecasting performance. Not one model stands out as best in all markets.}
\label{fig:maebar}
\end{figure}


\begin{table}[!htb]
\centering
\caption{Mean absolute error on log returns by model and metro (test period: 2023-01 to 2026-03). Bold indicates the best performance for each metro in the row.}
\label{tab:logmae}
\begin{tabular}{lrrrr}
\hline
Metro & OLS & ARIMA(AIC) & XGBoost & MLP(16,8) \\
\hline
San Francisco, CA & \textbf{0.00339} & 0.00497 & 0.00368 & 0.00455 \\
Austin, TX        & 0.00555 & 0.00518 & 0.00748 & \textbf{0.00488} \\
Cleveland, OH     & 0.00339 & \textbf{0.00226} & 0.00253 & 0.00487 \\
\hline
\end{tabular}
\end{table}


Figure~\ref{fig:preds} plots the estimated log returns for all four models compared to the actual time series in the test period.

\begin{figure}[!htb]
\centering
\includegraphics[width=0.68\textwidth]{predictions_vs_actual.png}
\caption{Estimated vs. Actual Monthly Log Returns of ZHVI 
in the Test Period (Jan. 2023 to Mar. 2026). The black line is the actual time series, while the colored lines represent model forecasts.}
\label{fig:preds}
\end{figure}


Our results show that no single model wins across all three markets. In San~Francisco, the pooled OLS model produces the best performance at \$3{,}898 per month, outperforming both non-linear models. In Austin, the MLP achieves the lowest dollar MAE at \$2{,}286 per month, with ARIMA a close second at \$2{,}371 per month --- a margin of less than \$100, so the MLP's edge over the pure time-series benchmark is small. In Cleveland, the MLP produces its worst performance, with a dollar MAE of \$1{,}064 per month, while ARIMA is the most convincing model for that metro at \$491 per month. Overall, these patterns demonstrate that model performance is geographically contingent.

\subsection{SHAP Feature Importance}

To add precision to our XGBoost results, we compute SHAP (SHapley Additive exPlanations) values over the test set and average contributions by metro. Three distinct patterns emerge. In \textbf{Cleveland}, \texttt{price\_cuts} carries the most weight (mean $|\text{SHAP}|=0.0039$), followed by \texttt{d\_log\_days\_pending} ($0.0011$), while housing starts are nearly irrelevant ($0.0002$); this is consistent with a friction-sensitive, supply-constrained market in which transaction signals dominate. In \textbf{Austin}, \texttt{price\_cuts} alone is roughly five times more influential than any other variable and accounts for much of the predictive advantage of listing-side data over purely macro-driven models, which fits Austin's volatile, demand-driven market where seller behavior is more informative than national macro series. In \textbf{San Francisco}, the metro indicator carries most of the weight ($0.0028$), followed by \texttt{price\_cuts} ($0.0018$) and \texttt{d\_mortgage30} ($0.0011$), suggesting that structural characteristics of the Bay Area (income distribution, price level) explain more of the monthly variation than any single time-varying indicator.

Table~\ref{tab:shap} lists the mean absolute SHAP values for the most important features per metropolitan area.

\begin{table}[!htb]
\centering
\caption{Mean absolute SHAP values by feature and metro 
(XGBoost, test set). Higher values indicate greater average 
contribution to the model's prediction.}
\label{tab:shap}
\begin{tabular}{lrrr}
\hline
Feature & San Francisco & Austin & Cleveland \\
\hline
\texttt{price\_cuts}           & 0.00180 & \textbf{0.00662} & \textbf{0.00391} \\
\texttt{metro\_San Francisco}  & \textbf{0.00279} & 0.00125 & 0.00122 \\
\texttt{d\_mortgage30}         & 0.00114 & 0.00088 & 0.00084 \\
\texttt{d\_log\_days\_pending} & 0.00086 & 0.00067 & 0.00110 \\
\texttt{d\_log\_new\_listings} & 0.00085 & 0.00063 & 0.00081 \\
\texttt{stlfsi}                & 0.00078 & 0.00071 & 0.00053 \\
\texttt{cpi\_yoy}              & 0.00046 & 0.00098 & 0.00053 \\
\texttt{d\_log\_houst}         & 0.00022 & 0.00020 & 0.00019 \\
\hline
\end{tabular}
\end{table}

\subsection{OLS Regression Table}

Table~\ref{tab:ols} reports the pooled OLS estimates ($N=168$, training sample 2018-01 to 2022-12). HAC (Newey--West, 3 lags) standard errors. The constant is omitted because the three metro dummies span the all-ones vector and serve as per-metro intercepts.

\begin{table}[!htb]
\centering
\small
\caption{Pooled OLS estimates. Dependent variable: monthly log return of ZHVI. ${}^{***}p<0.01$, ${}^{*}p<0.10$.}
\label{tab:ols}
\begin{tabular}{lrrr}
\hline
Predictor & Coef. & Std. Err. & $p$-value \\
\hline
\multicolumn{4}{l}{\textit{Time-varying predictors (lagged one month except CPI YoY)}} \\
\texttt{d\_mortgage30}        &  $0.0081^{***}$  & 0.002 & 0.000 \\
\texttt{d\_unrate\_nat}       &  $0.0005$        & 0.001 & 0.573 \\
\texttt{d\_unrate\_metro}     & $-0.0009$        & 0.001 & 0.347 \\
\texttt{d\_log\_houst}        & $-0.0068$        & 0.004 & 0.131 \\
\texttt{d\_log\_new\_listings}& $-0.0099^{*}$    & 0.005 & 0.062 \\
\texttt{d\_log\_days\_pending}&  $0.0032$        & 0.006 & 0.584 \\
\texttt{price\_cuts}          & $-0.1647^{***}$  & 0.015 & 0.000 \\
\texttt{stlfsi}               & $-0.0021^{***}$  & 0.001 & 0.006 \\
\texttt{cpi\_yoy}             & $-0.1010^{***}$  & 0.032 & 0.001 \\
\hline
\multicolumn{4}{l}{\textit{Metro fixed effects (per-metro intercepts)}} \\
\texttt{metro\_Austin}        &  $0.0439^{***}$  & 0.004 & 0.000 \\
\texttt{metro\_Cleveland}     &  $0.0380^{***}$  & 0.003 & 0.000 \\
\texttt{metro\_San~Francisco} &  $0.0296^{***}$  & 0.003 & 0.000 \\
\hline
$N$ & \multicolumn{3}{l}{168} \\
$R^{2}$ & \multicolumn{3}{l}{0.746} \\
Adj.~$R^{2}$ & \multicolumn{3}{l}{0.728} \\
Durbin--Watson & \multicolumn{3}{l}{0.637} \\
\hline
\end{tabular}
\end{table}

\subsection{Model Interpretability}

The four models differ in how much they support an economic reading of the predictions. \textbf{ARIMA} contains no exogenous regressors: its AR and MA coefficients describe the internal time-series structure of returns but cannot speak to which economic variables drive prices. The \textbf{MLP} is non-interpretable by design --- with ReLU activations and L2 regularization, individual weights have no economic meaning without post-hoc attribution, which we did not apply. \textbf{XGBoost} is interpretable in the sense we want via the SHAP attribution reported above (Table~\ref{tab:shap}). \textbf{OLS} is the only model whose coefficients are in principle interpretable as the marginal effect of each predictor on the log return. However, several diagnostics suggest reading them as descriptive partial correlations rather than causal estimates:

\begin{itemize}
    \item \textbf{Wrong-signed mortgage rate.} The estimated coefficient is positive ($\hat{\beta}_{\text{mortgage}} = +0.0081$), the opposite of what housing-demand theory predicts; this is most plausibly a confounding artefact of the post-COVID regime in which rates and prices both rose.
    \item \textbf{Statistical caveats.} Three regressors are borderline non-stationary by ADF (\texttt{d\_mortgage30} $p=0.063$, \texttt{d\_log\_days\_pending} $p=0.076$, \texttt{price\_cuts} $p=0.055$); the Durbin--Watson statistic of $0.64$ indicates strong positive residual autocorrelation; and national and metro unemployment have variance inflation factors near 8. Together these mean coefficients are inefficient and standard errors are too tight even with HAC correction.
    \item \textbf{Pooled across metros.} Each $\hat{\beta}$ is the average effect across San Francisco, Austin, and Cleveland; the pooled estimate masks any slope heterogeneity (this shows up indirectly through the per-metro intercepts but not through metro-specific slopes).
\end{itemize}
